\subsection{量化理论的公理}

量化理论是任何这样的理论 \(\mathfrak{C}\) ：其中方案S1至S4（§3，第1号）以及下面的方案S5提供隐式公理。

S5. 如果R是C中的一个关系，T是C中的一个项，且x是一个字母，则关系 \((T|x)R \Longrightarrow (\exists x)R\) 是一条公理。

这条规则确实是一个方案。因为设A是通过应用S5从C得到的一条公理；于是C中存在一个关系R、C中的一个项T以及一个字母x，使得A为 \((T|x)R \Longrightarrow (\exists x)R\) 。设U是C中的一个项，y是一个字母。我们将证明 \((U|y)A\) 也可以通过应用S5得到。利用CS1（§1，第2节）和CS8（第1节），我们可以仅限于x不同于y且不出现在U中的情形。设 \(R'\) 为关系 \((U|y)R\) ，并且令 \(T'\) 为项 \((U|y)T\) 。准则CS2（§1，第2条）和CS9（第1条）表明 \((U|y)A\) 与 \((T'|x)R' \Longrightarrow (\exists x)R'\) 相同。

方案S5表明，如果存在一个对象T，使得关系R（被视为表达x的一个性质）为真，那么R对于对象 \(\tau_{x}(R)\) 为真；这与我们赋予 \(\tau_{x}(R)\) 的直观含义相一致（§1，第3号，注记）。

